\section{The Gaussian Explicit Formula in a Fixed Fourier Convention}
\label{v4-arith:w5-i:weil-gaussians}

The zero test determines the time test by Fourier inversion.
For the two normalizations, put
\[
 \widehat f_{2\pi}(u)=\int f(x)e^{-2\pi iux}dx,\qquad
 (Sf)(t)=\frac1{2\pi}f\left(-\frac{t}{2\pi}\right).
\]
Then $\mathcal F_+(Sf)=\widehat f_{2\pi}$, where
$\mathcal F_+\varphi(u)=\int\varphi(t)e^{iut}dt$.
A change of variables gives $S(f*g)=Sf*Sg$ and $S(f^*)=(Sf)^*$.
It also gives $\|Sf\|_2^2=(2\pi)^{-1}\|f\|_2^2$.
Thus $S$ preserves the convolution algebra but is not a unitary
normalization of the scalar product.

\begin{definition}[the Gaussian pair]
\label{v4-arith:w5-i:def:gaussian-family}
For $\lambda>0$, define
\[
 f_\lambda(x)=e^{-\pi x^2/\lambda},\qquad
 h_\lambda(u)=\sqrt\lambda e^{-\pi\lambda u^2},\qquad
 \varphi_\lambda(t)=\frac1{2\pi}e^{-t^2/(4\pi\lambda)}.
\]
Then $h_\lambda=\widehat f_{\lambda,2\pi}
=\mathcal F_+\varphi_\lambda$.
All time tests belong to the rapid exponential space $\mathscr G$
of Proposition~\ref{v4-arith:sc:gaussian-extension}.
The symbol $W$ denotes the zero-side Weil distribution, so
$W(\varphi)=\sum_\rho\mathcal F_+\varphi(z_\rho)$.
\end{definition}

\begin{theorem}[the complete Gaussian identity]
\label{v4-arith:w5-i:thm:weil-gauss-closed}
With $z_\rho=-i(\rho-1/2)$, counted with multiplicity, one has
\begin{align}
 \sqrt\lambda\sum_\rho e^{-\pi\lambda z_\rho^2}
 ={}&2\sqrt\lambda e^{\pi\lambda/4}
 -\frac{\sqrt\lambda}{2\pi}\int_{\mathbb R}
 e^{-\pi\lambda u^2}k(u)du\notag\\
 &-\frac1\pi\sum_{n\ge2}\frac{\Lambda(n)}{\sqrt n}
              e^{-(\log n)^2/(4\pi\lambda)},
 \label{v4-arith:w5-i:eq:weil-id-gauss}
\end{align}
where $k(u)=\log\pi-\Re\psi_\Gamma(1/4+iu/2)$.
Every sum and integral converges absolutely, locally uniformly in
$\lambda>0$, together with every derivative in $\lambda$.
\end{theorem}
\begin{proof}
Completing the square gives the Fourier pair above.
The polar evaluations are $h_\lambda(\pm i/2)
=\sqrt\lambda e^{\pi\lambda/4}$.
The prime term is $\sum\Lambda(n)n^{-1/2}
(\varphi_\lambda(\log n)+\varphi_\lambda(-\log n))$,
which gives its exact coefficient and exponent.
The gamma term is $(2\pi)^{-1}\int h_\lambda k$.
Insert these three evaluations into the explicit formula on
$\mathscr G$.

On any interval $0<a\le\lambda\le b$, the zero terms and all
parameter derivatives are bounded by a polynomial in $|z_\rho|$
times $e^{-\pi a(\Re z_\rho)^2}$.
The bounded imaginary strip and the zero count make this summable.
For the prime terms, Gaussian decay in $\log n$ dominates every
power of $n^{-1}$, uniformly on $[a,b]$. Logarithmic growth of $k$
controls the gamma integral and its parameter derivatives.
\end{proof}

\begin{proposition}[the gamma concentration limit]
As $\lambda\to\infty$, the Guinand gamma term satisfies
\[
 -\frac{\sqrt\lambda}{2\pi}\int e^{-\pi\lambda u^2}k(u)du
 \longrightarrow-\frac{k(0)}{2\pi}<0.
\]
\end{proposition}
\begin{proof}
Put $v=\sqrt\lambda u$. Then $\int e^{-\pi v^2}dv=1$,
and $k(v/\sqrt\lambda)\to k(0)$.
For $\lambda\ge1$, logarithmic growth gives the common integrable
majorant $C e^{-\pi v^2}(1+\log(2+|v|))$.
Dominated convergence proves the limit.
\end{proof}

\begin{proposition}[a different finite part]
Define
\[
 A_\Gamma(\lambda)=\lim_{R\to\infty}
 \left\{\int_0^R(e^{-\pi x^2/\lambda}-1)
              \Im\psi_\Gamma(1/2+ix)dx+\frac\pi2 R\right\}.
\]
This limit exists and equals
\[
 A_\Gamma(\lambda)=\frac\pi2\int_0^\infty
                  e^{-\pi x^2/\lambda}\tanh(\pi x)dx
                  +\frac12\log2.
\]
It satisfies $A_\Gamma(\lambda)\sim\pi\sqrt\lambda/4$ as
$\lambda\to\infty$. Consequently $-\log\pi+A_\Gamma(\lambda)$
is not the gamma term in \eqref{v4-arith:w5-i:eq:weil-id-gauss}.
\end{proposition}
\begin{proof}
Digamma reflection at $1/2+ix$ gives
$\Im\psi_\Gamma(1/2+ix)=(\pi/2)\tanh(\pi x)$.
Separate the Gaussian integral from
$(\pi/2)\int_0^R(1-\tanh(\pi x))dx$.
The latter converges to $(\log2)/2$.
After $x=\sqrt\lambda v$, dominated convergence gives
$\lambda^{-1/2}\int_0^\infty e^{-\pi x^2/\lambda}
\tanh(\pi x)dx\to\int_0^\infty e^{-\pi v^2}dv=1/2$.
The preceding finite gamma limit proves the final assertion.
\end{proof}

\begin{proposition}[convolution and the full quadratic form]
For $\lambda,\mu>0$,
\[
 f_\lambda*f_\mu=\sqrt{\frac{\lambda\mu}{\lambda+\mu}}
 f_{\lambda+\mu},\qquad
 \varphi_\lambda*\varphi_\mu
 =\sqrt{\frac{\lambda\mu}{\lambda+\mu}}\varphi_{\lambda+\mu}.
\]
The corresponding Weil pairing is
\[
 W(\varphi_\lambda*\varphi_\mu)
 =\sqrt{\lambda\mu}\sum_\rho e^{-\pi(\lambda+\mu)z_\rho^2}.
\]
This is a zero sum, not the identically zero discrepancy between the
two sides of the explicit formula.
\end{proposition}
\begin{proof}
Complete the square in the time convolution. Apply the convolution
identity for $S$, then apply the Fourier product formula and the
absolutely convergent zero evaluation.
\end{proof}
